

\documentclass[12pt, a4paper, oneside]{report}

\usepackage{preamble}
\usepackage{graphs}
\usepackage{scripts}
\usepackage{alone}

 
\hypersetup{
    pdftitle   = {Anaerobic Co-digestion of Food Waste and Glycerol in Single-Stage System},
    pdfauthor  = {Diogenes Theodoros Bussador Matrakas},
    pdfsubject = {},
    pdfkeywords= {anaerobic digestion, biogas, crude glycerol, food waste, methane},
}

\begin{document}

\pagenumbering{roman}

\begin{titlepage}
    \centering

\begin{minipage}{0.45\textwidth}
        \raggedright
        \includegraphics[height=60pt]{./assets/logos/IPB_logo.png}
    \end{minipage}
    \hfill
    \begin{minipage}{0.45\textwidth}
        \raggedleft
        \includegraphics[height=60pt]{./assets/logos/UTFPR_logo.png}
    \end{minipage}

    \vspace{2.5cm}

{\LARGE \bfseries Anaerobic Co-digestion of Food Waste and Glycerol in Single-Stage System \par}

    \vspace{1.5cm}

{\Large \bfseries Diogenes Theodoros Bussador Matrakas \par}

    \vspace{2cm}

\large
    \textit{Dissertation submitted to Escola Superior Agrária de Bragança to obtain the Degree of Master in Biotechnological Engineering{} under the scope of the double diploma with Universidade Tecnológica Federal do Paraná.}

    \vspace{1.5cm}

Supervised by\\
    \textbf{Dr. Ramiro José Espinheira Martins}\\
    \textbf{Dr. Francisco Menino Destéfanis Vítola}

    \vfill

\normalsize
    {\fontsize{11}{13.2}\selectfont This dissertation does not include the comments and suggestions mentioned by the Jury.}

    \vspace{1cm}

\large
    \textbf{Bragança}\\
    \textbf{\monthname~\the\year}

    \vspace{0.5cm}
\end{titlepage}
 \cleardoublepage

\cleardoublepage

\cleardoublepage

\cleardoublepage

\glstocfalse
\printunsrtglossary[type=abbreviations, title={List of Abbreviations}]
\glstoctrue

\clearpage{}\chapter*{Acknowledgements}

\clearpage{}
\clearpage{}\chapter*{Abstract}

The growing generation of food waste and the surplus of crude glycerol from biodiesel production represent two significant organic waste streams requiring sustainable management solutions. This study evaluated the technical feasibility of a single-stage anaerobic co-digestion system treating food waste and glycerol as co-substrates, with the objective of verifing biogas production while providing a valorisation route for crude glycerol. The investigation was conducted in a 1 L continuous stirred-tank reactor operated under mesophilic conditions (35 °C) with a hydraulic retention time of 30 days. Food waste was collected from a university canteen, and crude glycerol was obtained from biodiesel production residues; anaerobic sludge from a municipal wastewater treatment plant served as inoculum. The reactor was operated in sequential phases: an initial adaptation period, co-digestion of food waste and glycerol, glycerol only feeding, and a post-eeding period. Biogas production was monitored continuously using an automatic methane potential test system, while process stability was assessed through measurements of pH, total alkalinity, volatile fatty acids, and their ratio. Substrate characterisation included total solids, volatile solids, density, total carbon, and total nitrogen analyses. Results demonstrated that co-digestion of food waste and glycerol yielded the highest methane production rates and cumulative volumes, with flow spikes immediately following feeding events indicating rapid substrate conversion. In contrast, glycerol-only feeding resulted in reduced biogas production despite continued methanogenic activity. Total alkalinity remained elevated following bicarbonate supplementation, maintaining pH around 7 throughout most of the experiment; volatile fatty acid concentrations fluctuated between 3000–5000 mg CH3COOH/L, with VFA/alkalinity ratios of 0.3–1.2. These findings confirm the hypothesis that glycerol can be useful for biogas production when co-digested with food waste or alone. The study demonstrates the technical viability of anaerobic co-digestion as a strategy for valorising crude glycerol within a circular bioeconomy framework, though further research is needed to optimise glycerol loading rates and develop robust feeding strategies for industrial application.

\vspace{1em}
\noindent\textbf{Keywords:} anaerobic digestion; biogas; crude glycerol; food waste; methane.
\clearpage{}

\begingroup
  \sisetup{
      output-decimal-marker={,},
      group-separator={.},
      group-minimum-digits=4,
      range-phrase={ a },
  }
  \clearpage{}\chapter*{Resumo}

A crescente geração de resíduos alimentares e o excedente de glicerol bruto proveniente da produção de biodiesel representam duas produções significativas de resíduos orgânicos que requerem soluções sustentáveis. Este estudo avaliou a viabilidade técnica de um sistema de codigestão anaeróbia de uma única etapa, tratando resíduos alimentares e glicerol como co-substratos, com o objetivo de verificar a produção de biogás e, simultaneamente, proporcionar uma via de valorização para o glicerol bruto. A investigação foi conduzida num reator contínuo de tanque agitado de 1 L, operado em condições mesofílicas (35 °C), com um tempo de retenção hidráulica de 30 dias. Os resíduos alimentares foram recolhidos numa cantina universitária e o glicerol bruto foi obtido a partir de resíduos da produção de biodiesel; lamas anaeróbias de uma estação de tratamento de águas residuais municipais serviram como inóculo. O reator foi operado em fases sequenciais: um período inicial de adaptação, codigestão de resíduos alimentares e glicerol, alimentação exclusiva com glicerol e um período pós-alimentação. A produção de biogás foi monitorizada continuamente através de um \textit{automatic methane potential test system}, enquanto a estabilidade do processo foi avaliada por medições de pH, alcalinidade total, ácidos graxos voláteis e respectiva razão. A caracterização dos substratos incluiu análises de sólidos totais, sólidos voláteis, massa volúmica, carbono total e nitrogênio total. Os resultados demonstraram que a codigestão de resíduos alimentares e glicerol originou as taxas de produção de metano e os volumes acumulados mais elevados, com picos imediatamente após as alimentações, indicando uma rápida conversão do substrato. Em contraste, a alimentação exclusiva com glicerol resultou numa produção reduzida de biogás, apesar da atividade metanogénica continuada. A alcalinidade total manteve-se elevada após a suplementação com bicarbonato, mantendo o pH em torno de 7 durante a maior parte da experiência; as concentrações de ácidos graxos voláteis flutuaram entre 3000–5000 mg CH3COOH/L, com razões AGV/alcalinidade de 0,3–1,2. Estes resultados confirmam a hipótese de que o glicerol pode ser útil para a produção de biogás quando codigerido com resíduos alimentares ou isoladamente. O estudo demonstra a viabilidade técnica da codigestão anaeróbia como estratégia de valorização do glicerol bruto no âmbito de uma bioeconomia circular, embora seja necessária investigação adicional para otimizar as cargas de glicerol e desenvolver estratégias de alimentação robustas para aplicação industrial.

\vspace{1em}
\noindent\textbf{Palavras-chave:} biogás; digestão anaeróbia; glicerol bruto; metano; resíduos alimentares.
\clearpage{}
\endgroup

\pagenumbering{arabic}
\setcounter{page}{1}

\clearpage{}\chapter{Introduction}

The generation of organic waste has become one of the defining environmental challenges of our time, particularly as urbanization and industrial activity continue to expand. A major contributor is FW, a global stream that in the European Union alone is estimated at 88 Mt/yr, excluding inedible parts removed from the supply chain \parencite{scherhaufer2018}. When considering only post-processing activities (wholesale, retail, food service and households) approximately 66.3 Mt of FW are produced annually in the EU-27+1 \parencite{albizzati2021}. The scale of this loss highlights the urgent need for improved management strategies that move beyond simple disposal.

The environmental burden of FW is substantial. It accounts for roughly 15–16 % of the total climate impact of the entire food supply chain, translating to about 186 Mt CO2-eq of greenhouse gas emissions, alongside significant contributions to acidification and eutrophication \parencite{scherhaufer2018}. The majority of these impacts arise during primary production, making prevention the single most effective strategy; however, when waste cannot be avoided, its valorisation through anaerobic digestion or composting is favoured over landfill disposal \parencite{albizzati2021, scherhaufer2018}. Despite this, much FW continues to be landfilled or incinerated, forfeiting the opportunity to recover its embedded energy and nutrients.

Parallel to the FW challenge, the rapid expansion of biodiesel production has generated an excess of another organic by-product: C-GL. Formed in the transesterification of triglycerides, approximately 10 kg of C-GL are produced for every 100 kg of biodiesel \parencite{garlapati2016}. Global biodiesel output, driven by renewable energy mandates, has caused the GL market to swell with C-GL supply projected to reach about 6.33 Mt by 2025, of which nearly 680 kt are expected to be highly impure due to the shift towards waste-based feedstocks \parencite{attarbachi2023}.

Crude glycerol composition varies widely depending on the oil source and catalyst used, but it typically contains methanol, soaps, free fatty acids, salts and other MONG compounds \parencite{anand2012, attarbachi2023}. These impurities are known to inhibit microbial growth and fermentation, as demonstrated by the severe reduction in 1,3-PDO production when untreated C-GL from jatropha, soybean or sunflower oils was used \parencite{anand2012}. Purification to technical or pharmaceutical grade is energy-intensive and often uneconomical for small and medium biodiesel plants \parencite{attarbachi2023}, prompting the search for direct, cost-effective valorisation routes. Because of its high organic content and ready biodegradability, C-GL has attracted attention as an ideal co-substrate for anaerobic digestion, where it can serve as a concentrated carbon source while being simultaneously detoxified by dilution and microbial synergy \parencite{gebreegziabher2025, tomczak2025}.

Anaerobic digestion is a mature biological process in which consortia of microorganisms convert organic matter into biogas, primarily CH4 (55–65 %) and CO2, under O2-free conditions, passing through the sequential stages of hydrolysis, acidogenesis, acetogenesis and methanogenesis \parencite{gebreegziabher2025}. Anaerobic digestion offers a unique combination of benefits: it produces renewable, storable energy that can replace fossil fuels; generates a nutrient-rich digestate that can be used as organic fertiliser; and reduces the environmental load of organic wastes by avoiding landfilling and its associated CH4 emissions \parencite{scherhaufer2018, gebreegziabher2025}.

Within this framework, anaerobic co-digestion, the simultaneous treatment of two or more substrates, has emerged as a particularly effective strategy. Co-digesting C-GL with nitrogen-rich materials such as animal manure, sewage sludge or FW balances the C/N, dilutes potential inhibitors, provides essential macro and micro nutrients, and fosters synergistic microbial interactions \parencite{tomczak2025, gebreegziabher2025}. The result can be a dramatic improvement in biogas and CH4 yield: studies report increases ranging from 14–226 % when C-GL is added to a variety of base substrates, with CH4 contents generally exceeding 60 % as long as GL does not exceed about 8 % (volume/volume) of the feed \parencite{gebreegziabher2025, tomczak2025}.

However, excessive GL loads can lead to VFA accumulation, pH drop and methanogen inhibition, underscoring the need for careful optimisation \parencite{gebreegziabher2025, tomczak2025}. Overall, the convergence of abundant FW and surplus C-GL supplies, together with the proven capability of AD to convert these streams into renewable CH4 and fertiliser, represents a tangible route towards a circular, low-carbon bioeconomy.

\chapter{Objectives}

\section{General Objectives}

The primary objective of this dissertation project is to evaluate the efficiency of a single-stage anaerobic co-digestion system using organic FW and C-GL as substrates. The investigation will be conducted by determining biogas production, with a focus on industrial feasibility.

The central hypothesis is that the controlled addition of C-GL, a by-product of the biodiesel industry, helps biogas production due to its high energy content and rapid biodegradability, provided that adequate operational conditions are maintained to avoid VFA overload.

\section{Specific Objectives}

\begin{enumerate}
    \item Characterize the substrates (FW and C-GL), evaluating physicochemical parameters relevant to anaerobic digestion.
    \item Operate the single-stage anaerobic co-digestion system, assessing process stability.
    \item Evaluate the techno-economic feasibility of using C-GL as a co-substrate, considering potential for industrial application.
\end{enumerate}

\clearpage{}
\clearpage{}\chapter{Literature Review}

The rapid expansion of biodiesel production has generated an oversupply of C-GL, a by-product of the transesterification output, this ?? shows the steps to produce biodiesel and when GL is producesd \parencite{garlapati2016}. This waste stream, loaded with impurities such as methanol, soaps, and hydroxides, the ?? shows the processes and what they contaminate the GL with, poses environmental disposal challenges, also with the variation of impurities depending on the source, there's no standard method to use this source of carbon \parencite{anand2012, gebreegziabher2025}. Anaerobic digestion aligns with circular economy principles by converting organic residues into biogas while recycling nutrients, making C-GL a promising co-substrate \parencite{gebreegziabher2025}.

\begin{figure}[htbp]
    \centering
    \resizebox{0.65\textwidth}{!}{% TIKZ-FAILED: assets/graphs/garlapati2016-flow-gl
    }
    \caption{crude glycerol from biodisel production (Adapted from: \textcite{garlapati2016})}
    
\end{figure}

\begin{figure}[htbp]
    \centering
    \resizebox{\textwidth}{!}{% TIKZ-FAILED: assets/graphs/attarbachi2023-flow-gl
    }
    \caption{Contaminations of the crude glycerol (Adapted from: \textcite{attarbachi2023})}
    
\end{figure}

\textcite{gebreegziabher2025} reported that co-digestion of C-GL can boost CH4 production compared to mono-digestion of the base substrate. \textcite{tomczak2025} reviewed one-stage AD experiments and concluded that a CH4 content above 60 % could be consistently obtained when the C-GL fraction did not exceed 8 % (volume/volume) of the feed, the flow chart by \textcite{li2011} on the ?? it shows that using AD there could be up to 70 % CH4 production. However, the same review emphasized that increasing the GL loading beyond a system specific threshold may provoke inhibitory effects, likely due to the accumulation of VFAs and the impurities present in C-GL \parencite{anand2012, tomczak2025}. The heterogeneous composition of C-GL, affected by oil feedstock, catalyst, and biodiesel processing, demands careful feedstock characterization and, if necessary, simple pretreatment to avoid process upsets \parencite{attarbachi2023}.

\begin{figure}[htbp]
    \centering
    \resizebox{0.65\textwidth}{!}{% TIKZ-FAILED: assets/graphs/li2011-flow-ferm
    }
    \caption{CH4 production from AD (Adapted from: \textcite{li2011})}
    
\end{figure}

Liquid AD, typically conducted in a CSTR, operates at total solid concentrations of 0.5–15 % and is generally used for treating sewage sludge, manure, and FW \parencite{li2011}. The CSTR’s continuous mixing provides uniform distribution of substrates, microorganisms, and heat, which is essential when feeding a fermentable substrate. Temperature is a critical operational parameter; the activity of methanogenic consortia is strongly temperature dependent, and both mesophilic and thermophilic regimes have been applied to AD \parencite{choorit2007}. \textcite{tomczak2025} noted that the majority of GL co-digestion studies have been performed under mesophilic conditions, while thermophilic operation, which could offer higher conversion rates, still requires more systematic investigation. The CSTR platform is well suited to explore such temperature effects because its design allows easy control of heating and hydraulic retention time.

The selection of a robust microbial inoculum is decisive for rapid start-up and stable performance of a CSTR fed with GL. Anaerobic sludge from municipal WWTPs is the most common source of inoculum because it harbours a diverse consortium of hydrolytic, acidogenic, and methanogenic microorganisms, the flow cahrd from \textcite{gebreegziabher2025} on the ?? shows each step of the AD until the CH4 production. \textcite{posmanik2020} demonstrated that granular sludge, obtained from up-flow anaerobic sludge blanket reactors treating wastewater, outperforms suspended biomass as an inoculum. Although their work focused on batch tests, the findings strongly suggest that using granular sludge from a WWTP as seed for a CSTR could accelerate acclimation, improve process stability, and reduce the risk of acidification during load increases.

\begin{figure}[htbp]
    \centering
    \resizebox{0.65\textwidth}{!}{% TIKZ-FAILED: assets/graphs/gebreegziabher2025-flow-ferm
    }
    \caption{Steps to produce CH4 from AD (Adapted from: \textcite{gebreegziabher2025})}
    
\end{figure}

When GL starts being degraded it makes the AD process susceptible to VFA accumulation and pH drop if the organic loading rate is not properly controlled. Therefore, close monitoring of process indicators is essential. \textcite{ribas2007} evaluated several methods for determining these parameters specifically in anaerobic reactors, highlighting the importance of accurate VFA and Alk measurements for timely detection of process imbalance. Using the method from DiLallo \& Albertson \parencite{ribas2007} provide a validated method for measuring VFAs and Alk. Maintaining a favourable vfa/alk is particularly critical when co-digesting GL, as excess VFAs can inhibit methanogenesis \parencite{tomczak2025}.

The available literature confirms that C-GL is a high-energy substrate that can significantly boost CH4 production in liquid AD systems. A CSTR provides the mixing and temperature control required to harness this potential. Using granular sludge from a WWTP as the inoculum can enhance start-up speed and methanogenic activity, as demonstrated by comparative studies. To unlock the full benefits while avoiding inhibition, the GL loading must be kept below documented thresholds, and the reactor state should be continuously tracked through standard measurements of VFAs and Alk. Further research is needed to optimise thermophilic CSTR operation and to establish long-term performance data for granular-sludge-inoculated reactors treating C-GL.
\clearpage{}
\clearpage{}\chapter{Materials and Methods}

\section{Collection and Preparation of the Food Waste}

Food waste was collected from the canteen at the IPB, prioritizing residues with high proportions of rice, pasta, and bread to maximize the sugar content and minimize the nitrogen content. The waste was diluted with approximately 2.5 L of distilled water and homogenised using a blender to obtain a liquid paste. It was then stored in a refrigerator at 4 °C to prevent spoilage and suppress microbial and fungal activity \parencite{degueurce2020}.

\section{Crude Glycerol Collection}

A 6 L bottle of C-GL, obtained from previous biodiesel production experiments at IPB, was provided for this study. C-GL from biodiesel typically contains impurities such as catalyst residues \parencite{attarbachi2023}, which may influence its behaviour in anaerobic digestion.

\section{Inoculum Collection}

The inoculum was collected from a anaerobic digestor from Águas do Norte wastewater treatment plant in Bragança. This ensured the availability of a microbial population for the anaerobic digestion \parencite{posmanik2020}.

\section{Physical and Chemical Characterisation}

The TS, VS, FS, and density were determined for the FW, C-GL, and inoculum. Total organic carbon and TN were determined for the FW and C-GL.

\subsection{TS and VS Determination}

Total solids and Volatile solids were measured using a sequential gravimetric method based on Standard Methods 2540 B and 2540 E \parencite{apha2023}. The method involves controlled drying and calcination steps to quantify TS and VS fractions.

Initially, porcelain crucibles were placed in a muffle furnace at 550 °C for 1 h. For cooling, they were transferred to an oven at 105 °C for 15 min and then moved to a desiccator until they reached room temperature; their empty masses were recorded. For each material, a sample was placed into three crucibles and the wet mass was measured. The crucibles were then dried in an oven at 105 °C for 2 h (C-GL required 6–8 h to reach constant mass). After drying, the crucibles were cooled in a desiccator to room temperature and weighed again to determine TS.

The same crucibles were subsequently placed in a muffle furnace at 550 °C for 50 min. Afterwards they were transferred to an oven at 105 °C for 15 min and then moved to a desiccator until they reached room temperature; the final masses were recorded to calculate VS.

Total solids, Volatile solids and Fixed solids were calculated as follows:
\[
    \mathrm{TS} = \frac{M_{105} - M_{\text{crucible}}}{V_{\text{sample}}} \quad \left[g/L\right]
\]
\[
    \mathrm{VS} = \frac{M_{105} - M_{550}}{V_{\text{sample}}} \quad \left[g/L\right]
\]
\[
    \mathrm{FS} = \frac{M_{550} - M_{\text{crucible}}}{V_{\text{sample}}} \quad \left[g/L\right]
\]
where $M_{\text{wet}}$ is the mass of crucible plus wet sample, $M_{105}$ is the mass after drying at 105 °C, $M_{550}$ is the mass after ignition at 550 °C, $M_{\text{crucible}}$ is the mass of the empty crucible and $V_{\text{sample}}$ is the volume of the sample. The volume is in L and the weight is in g. The results are on the ??.

\begingroup
\tablesetup
\footnotesize
\setlength{\tabcolsep}{3pt}

    
        \begin{tabular}{llllllllllll}

        \caption{Total solids, Volatile solids and Fixed solids values}
         \\

        \toprule
        \multicolumn{2}{c}{\thead{Sample}} &
        \thead{$M_{\text{crucible}}$\\(g)} &
        \thead{$M_{\text{wet}}$\\(g)} &
        \thead{$M_{\text{sample}}$\\(g)} &
        \thead{$V_{\text{sample}}$\\(L)} &
        \thead{$M_{105}$\\(g)} &
        \thead{TS\\(g/L)} &
        \thead{Average\\TS\\(g/L)} &
        \thead{$M_{550}$\\(g)} &
        \thead{VS\\(g/L)} &
        \thead{Average\\VS\\(g/L)} \\
        \midrule
        \endfirsthead

        \multicolumn{12}{l}{\textit{?? continued}} \\
        \toprule
        \multicolumn{2}{c}{\thead{Sample}} &
        \thead{$M_{\text{crucible}}$\\(g)} &
        \thead{$M_{\text{wet}}$\\(g)} &
        \thead{$M_{\text{sample}}$\\(g)} &
        \thead{$V_{\text{sample}}$\\(L)} &
        \thead{$M_{105}$\\(g)} &
        \thead{TS\\(g/L)} &
        \thead{Average\\TS\\(g/L)} &
        \thead{$M_{550}$\\(g)} &
        \thead{VS\\(g/L)} &
        \thead{Average\\VS\\(g/L)} \\
        \midrule
        \endhead

        \bottomrule
        \endlastfoot

        \mbox{Sample} & \mbox{} & \mbox{m0 (g)} & \mbox{m1 (g)} & \mbox{m sub (g)} & \mbox{V sub (L)} & \mbox{m2 (g)} & \mbox{TS (g/L)} & \mbox{media TS (g/L)} & \mbox{m3 (g)} & \mbox{VS (g/L)} & \mbox{media VS (g/L)} \\
\mbox{Food Waste} & \mbox{1} & \mbox{78.731} & \mbox{105.9875} & \mbox{27.2565} & \mbox{0.0255} & \mbox{91.0021} & \mbox{481.0194} & \mbox{536.3998} & \mbox{79.2024} & \mbox{462.5408} & \mbox{520.8068} \\
\mbox{} & \mbox{2} & \mbox{84.3206} & \mbox{113.2979} & \mbox{28.9773} & \mbox{0.0271} & \mbox{98.8684} & \mbox{536.3998} & \mbox{} & \mbox{84.7435} & \mbox{520.8068} & \mbox{} \\
\mbox{} & \mbox{3} & \mbox{77.5887} & \mbox{108.8399} & \mbox{31.2512} & \mbox{0.0292} & \mbox{93.4417} & \mbox{541.9933} & \mbox{} & \mbox{78.1848} & \mbox{521.6134} & \mbox{} \\
\mbox{Crude Glycerol} & \mbox{1} & \mbox{84.2193} & \mbox{88.6434} & \mbox{4.4241} & \mbox{0.005} & \mbox{88.4966} & \mbox{863.4402} & \mbox{867.724} & \mbox{84.2144} & \mbox{864.4293} & \mbox{868.6548} \\
\mbox{} & \mbox{2} & \mbox{73.478} & \mbox{77.9874} & \mbox{4.5094} & \mbox{0.005} & \mbox{77.8594} & \mbox{867.724} & \mbox{} & \mbox{73.4733} & \mbox{868.6548} & \mbox{} \\
\mbox{} & \mbox{3} & \mbox{86.2012} & \mbox{90.6107} & \mbox{4.4095} & \mbox{0.0049} & \mbox{90.5034} & \mbox{871.3421} & \mbox{} & \mbox{86.1958} & \mbox{872.4358} & \mbox{} \\
\mbox{Inoculum} & \mbox{1} & \mbox{78.731} & \mbox{105.9875} & \mbox{27.2565} & \mbox{0.0255} & \mbox{91.0021} & \mbox{481.0194} & \mbox{536.3998} & \mbox{79.2024} & \mbox{462.5408} & \mbox{520.8068} \\
\mbox{} & \mbox{2} & \mbox{84.3206} & \mbox{113.2979} & \mbox{28.9773} & \mbox{0.0271} & \mbox{98.8684} & \mbox{536.3998} & \mbox{} & \mbox{84.7435} & \mbox{520.8068} & \mbox{} \\
\mbox{} & \mbox{3} & \mbox{77.5887} & \mbox{108.8399} & \mbox{31.2512} & \mbox{0.0292} & \mbox{93.4417} & \mbox{541.9933} & \mbox{} & \mbox{78.1848} & \mbox{521.6134} & \mbox{} \\

        \end{tabular}
    

\endgroup

\subsection{Density Determination}

Density was determined by a gravimetric volumetric method using graduated cylinders. Dry cylinders were weighed empty, then filled with distilled water at 20 °C and weighed again to verify the reading accuracy (density of water at $\rho_{20\,°C} = 998.21 g/L$). After drying, each cylinder was filled with a known volume of sample and weighed. The density of each material was calculated using
\[
    \rho = \frac{M_{\text{total}} - M_{\text{cylinder}}}{V_{\text{sample}}} \quad \left[g/L\right]
\]
where $M_{\text{total}}$ is the mass of cylinder plus sample, $M_{\text{cylinder}}$ is the empty cylinder mass, and $V_{\text{sample}}$ is the measured volume, the masses were messured in g and the volume in L. The results are on the ??.

\begingroup
\tablesetup
\footnotesize
\setlength{\tabcolsep}{3pt}

    \begin{tabular}{lllllll}
        \caption{Density values}
         \\

        \toprule
        \multicolumn{2}{c}{\thead{Sample}} &
        \thead{$M_{\text{cylinder}}$ (g)} &
        \thead{$M_{\text{total}}$ (g)} &
        \thead{$V_{\text{sample}}$ (L)} &
        \thead{Density\\(g/L)} &
        \thead{Average\\Density\\(g/L)} \\
        \midrule
        \endfirsthead

        \multicolumn{7}{l}{\textit{?? continued}} \\
        \toprule
        \multicolumn{2}{c}{\thead{Sample}} &
        \thead{$M_{\text{cylinder}}$ (g)} &
        \thead{$M_{\text{total}}$ (g)} &
        \thead{$V_{\text{sample}}$ (L)} &
        \thead{Density\\(g/L)} &
        \thead{Average\\Density\\(g/L)} \\
        \midrule
        \endhead

        \bottomrule
        \endlastfoot

        \mbox{Sample} & \mbox{} & \mbox{m0 (g)} & \mbox{m2 (g)} & \mbox{V (L)} & \mbox{Density (g/L)} & \mbox{Avarage Density (g/L)} \\
\mbox{Food Waste} & \mbox{1} & \mbox{136} & \mbox{242} & \mbox{0.0992} & \mbox{1068.7905} & \mbox{1068.4376} \\
\mbox{} & \mbox{2} & \mbox{116} & \mbox{223} & \mbox{0.1002} & \mbox{1068.0847} & \mbox{} \\
\mbox{Crude Glycerol} & \mbox{1} & \mbox{114.21} & \mbox{203.47} & \mbox{0.0999} & \mbox{893.2353} & \mbox{893.074} \\
\mbox{} & \mbox{2} & \mbox{115.74} & \mbox{204.44} & \mbox{0.0993} & \mbox{892.9127} & \mbox{} \\
\mbox{Inoculum} & \mbox{1} & \mbox{136} & \mbox{242} & \mbox{0.0992} & \mbox{1068.7905} & \mbox{1068.4376} \\
\mbox{} & \mbox{2} & \mbox{116} & \mbox{223} & \mbox{0.1002} & \mbox{1068.0847} & \mbox{} \\

    \end{tabular}

\endgroup

\subsection{TOC and TN Determination}

Total organic carbon, TC and TN were measured using a Shimadzu TOC-5000A analyser. Prior to analysis, the FW was diluted 400 X and the C-GL was diluted 800 X to bring the concentrations within the linear range of the instrument. The results are on the ??.

\begingroup
\tablesetup
\footnotesize
\setlength{\tabcolsep}{3pt}

    \begin{tabular}{llllllllllllll}

        \caption{Total carbon and Total nitrogen values}
         \\

        \toprule
        \multicolumn{2}{c}{\thead{Sample}} &
        \thead{TOC\\(mg/L)} &
        \thead{TC\\(mg/L)} &
        \thead{Average\\TC\\(mg/L)} &
        \thead{TN\\(mg/L)} &
        \thead{Average\\TN\\(mg/L)} \\
        \midrule
        \endfirsthead

        \multicolumn{7}{l}{\textit{?? continued}} \\
        \toprule
        \multicolumn{2}{c}{\thead{Sample}} &
        \thead{TOC\\(mg/L)} &
        \thead{TC\\(mg/L)} &
        \thead{Average\\TC\\(mg/L)} &
        \thead{TN\\(mg/L)} &
        \thead{Average\\TN\\(mg/L)} \\
        \midrule
        \endhead

        \bottomrule
        \endlastfoot

        \mbox{Sample} & \mbox{} & \mbox{TOC (mg/L)} & \mbox{TC (mg/L)} & \mbox{Avarage TC (mg/L)} & \mbox{TN (mg/L)} & \mbox{Avarage TN (mg/L)} \\
\mbox{Food Waste} & \mbox{1} & \mbox{5.039} & \mbox{2017.6} & \mbox{1778.6} & \mbox{275.64} & \mbox{281.76} \\
\mbox{} & \mbox{2} & \mbox{3.927} & \mbox{1539.6} & \mbox{} & \mbox{287.88} & \mbox{} \\
\mbox{Crude Glycerol} & \mbox{1} & \mbox{8.808} & \mbox{7049.6} & \mbox{5523.2} & \mbox{470} & \mbox{488.6} \\
\mbox{} & \mbox{2} & \mbox{4.992} & \mbox{3996.8} & \mbox{} & \mbox{507.2} & \mbox{} \\
    \end{tabular}

\endgroup

\section{Experiment Configuration}

The reactor had a useful volume of 1 L and was operated as a CSTR under anaerobic conditions. Mixing was provided by a magnetic stir bar, and the overflow exit tube was kept closed with a valve to maintain an anaerobic environment. As you can see in this ??:

\begin{figure}[H]
    \centering
    \resizebox{\textwidth}{!}{% TIKZ-FAILED: assets/graphs/reactor-model
    }
    \caption{Schematic of the experimental setup used}
    
\end{figure}

The reactor was kept at 35 °C, close to the optimal temperature for mesophilic microorganisms \parencite{choorit2007}. The OLR was set with a max of 4 kg VS/m³/d, with a HRT of 30 days, conditions for stable co-digestion of FW and GL \parencite{tomczak2025}. The pH was maintained around 7 by incorporating KHCO3 into the feed solution at a concentration of 10 g/L, and adding KHCO3 to the reactor whenever the pH dropped lower than 6.

The OLR and HRT:

\[
    \mathrm{HRT} = \frac{V}{Q} \quad \left[d\right]
\]
\[
    \mathrm{OLR} = \frac{S_{\text{in}} \times Q}{V} \quad \left[kg VS/m³/d\right]
\]
where $S_{\text{in}}$ is the influent VS concentration (kg/m³), $Q$ is the volumetric flow rate (m³/d), and $V$ is the reactor working volume (m³).

\subsection{Feeding Solution}

The feed was designed to achieve a C/N ratio of 25:1, which is optimal for AD \parencite{li2011}. Characterisation of the FW and C-GL showed that both substrates contained nitrogen, making it difficult to balance the C/N ratio using both of them together. Therefore, P-GL was used as an additive to allow precise control of the C/N. Based on the chosen HRT and the daily volumetric flow rate, the reactor was fed once per day to ensure a flow rate high enough to transport the solid FW suspension through the feeding tube.

The feed was made with a restriction of needing to have less than 15 % of TS%. The amount of OLR and if KHCO3 was added, was decided by the pH and how was the CH4 production of the reactor. If it was too acidic, it was used a lower OLR with KHCO3, if it was at the right amount it was keep the same amount, and if it was too basic it was added a higher OLR without KHCO3. The feed storage always had some remainig feed, and the new one was added on top, diluting or concentrating the feed going to the reactor.

The feed was also suplemented with nutrients to complement the feed and create a better feed for the microorganisms. ?? shows the nutrients used and there ratios.

\begingroup
\tablesetup
\footnotesize
\setlength{\tabcolsep}{3pt}

    \begin{tabular}{ll}

        \caption{Nutrients used and there ratios, for a 25:1 C/N. (Adapted from: \textcite{safaa2025})}
         \\

        \toprule
        \thead{Nutrients} &
        \thead{Mass (mg/L)} \\
        \midrule
        \endfirsthead

        \multicolumn{2}{l}{\textit{?? continued}} \\
        \toprule
        \thead{Nutrients} &
        \thead{Mass (mg/L)} \\
        \midrule
        \endhead

        \bottomrule
        \endlastfoot

        Nutrients & Mass (mg/L) \\
MgSO4.7H2O & 54 \\
FeSO4.7H2O & 1493.7143 \\
CoCl2.6H2O & 1.2 \\
ZnCl2 & 30 \\
NiCl2.6H2O & 0.03 \\
EDTA & 0.6 \\
CuSO4.5H2O & 26.3714 \\

    \end{tabular}

\endgroup

\subsection{Biogas Quantification}

The AMPTS II system (hardware~??, software~??), was provided to meassure the gas flow that was comming out of the reactor, with a CO2 fixing unit using a solution of 3 MNa OH.

\begin{figure}[H]
    \centering
    \begin{subfigure}[t]{0.45\textwidth}
        \centering
        \includegraphics[width=\linewidth]{./assets/photos/ampts/ampts.jpg}
        \caption{AMPTS II gas flow messurer}
        
    \end{subfigure}\hspace{1em}
    \begin{subfigure}[t]{0.45\textwidth}
        \centering
        \includegraphics[width=\linewidth, height=0.35\textheight, keepaspectratio]{./assets/photos/ampts/ampts_co2.jpg}
        \caption{AMPTS II CO2 trap}
        
    \end{subfigure}
    \caption{AMPTS II hardware used}
    
\end{figure}

\begin{figure}
    \centering
    \includegraphics[width=\linewidth]{./assets/photos/ampts/ampts_sys-print.png}
    \caption{AMPTS II software used}
    
\end{figure}

\subsection{Analysis Methods}

It was used the method of DiLallo \& Albertson \parencite{ribas2007}, to measure the Alk and VFA. The sample collected was first centrifuged, at 2604 G for 30 min, to filter the particulated before doing the measurements.

\[
    \mathrm{TA} = \frac{N_\text{acid} * V_\text{acid} * 50000}{V_\text{sample}} \quad \left[mg Ca CO3/L\right]
\]
\[
    \mathrm{VFA} = \frac{N_\text{base} * V_\text{base} * 60000}{V_\text{sample}} \quad \left[mg CH3COOH/L\right]
\]

The volumes were measured in m L, the sample volume was always 50 m L, the normality (N) of the acid and base, at first was 0.1, but then it was changed to 1. The ?? is expressed in mg Ca CO3/L, and the ?? is expressed in mg CH3COOH/L.

\subsection{Subestrate used as feed}

?? shows all substrates used to feed the reactor and when they were made.

\begingroup
\sisetup{}
\footnotesize
\setlength{\tabcolsep}{3pt}

    \begin{tabular}{lllllll}

    \caption{Description of the feed used throughout the experimental period}
     \\

    \toprule
    \thead{Day} &
    \thead{Substrate 1} &
    \thead{Subestrate 2} &
    \thead{Substrate 3} &
    \thead{Extra\\ Addition} &
    \thead{OLR\\ \& TS\%} &
    \thead{Volume\\ Made} \\
    \midrule
    \endfirsthead

    \multicolumn{6}{l}{\textit{?? continued}} \\
    \toprule
    \thead{Day} &
    \thead{Substrate 1} &
    \thead{Subestrate 2} &
    \thead{Substrate 3} &
    \thead{Extra\\ Addition} &
    \thead{OLR\\ \& TS\%} &
    \thead{Volume\\ Made} \\
    \midrule
    \endhead

    \bottomrule
    \endlastfoot

    0(a) &
    FW, 156 g &
    C6H12O6, 1.9 g &
    &
    Ca(HCO3)2, 5 g &
    1.5 \& 15 &
    500 m L \\
    
    4 &
    C-GL, 23.2 m L &
    P-GL(b), 31.5 m L &
    &
    KHCO3, 1.9 g &
    0.8 \& 10 &
    250 m L \\
    
    11 &
    FW, 12.54 g &
    C-GL, 9.02 m L &
    P-GL, 2.6 m L &
    KHCO3, 1.57 g &
    0.13–0.4 \& 5–10(c) &
    250 m L \\
    
    18 &
    FW, 10.88 g &
    C-GL, 2.89 m L &
    P-GL, 1.44 m L &
    &
    0.0256 \& 2 &
    250 m L \\
    
    25 &
    FW, 31.12 g &
    C-GL, 44 m L &
    P-GL, 29 m L &
    &
    1.4 \& 10 &
    250 m L \\
    
    32 &
    FW, 31.01 g &
    C-GL, 44 m L &
    P-GL, 28.9 m L &
    &
    1.4 \& 10 &
    250 m L \\
    
    39 &
    FW, 46.52 g &
    C-GL, 66 m L &
    P-GL, 43 m L &
    KHCO3, 2.6 g &
    3.1 \& 15 &
    250 m L \\
    
    46 &
    FW, 50.07 g &
    C-GL, 66 m L &
    P-GL, 43 m L &
    KHCO3, 2.67 g &
    3.16 \& 15 &
    250 m L \\
    
    53(d) &
    C-GL, 55 m L &
    P-GL, 11.5 m L &
    &
    KHCO3, 5.38 g &
    0.42 \& 10 &
    500 m L \\
    
    60 &
    C-GL, 55 m L &
    P-GL, 12 m L &
    &
    KHCO3, 5.18 g &
    0.42 \& 10 &
    500 m L(e) \\
    
    67 &
    C-GL, 41 m L &
    P-GL, 9 m L &
    &
    KHCO3, 2.56 g &
    0.95 \& 15 &
    250 m L \\
    
    74 &
    C-GL, 60 m L &
    P-GL, 12.5 m L &
    &
    KHCO3, 2.57 g &
    1.37 \& 15 &
    250 m L \\

    \end{tabular}

    

\endgroup
\clearpage{}
\clearpage{}\chapter{Results}

\section{Reactor Condition and Visual Observations}

The reactor was inoculated with anaerobic sludge and initially contained a murky, light brown suspension (??). Within three days of operation, the colour changed to black (??), which is characteristic of active anaerobic digestion and is typically caused by the precipitation of metal sulphides, particularly iron sulphide, under reducing conditions. This visual change indicated the establishment of a suitable anaerobic environment and the onset of methanogenic activity.

\begin{figure}[H]
    \centering
    \begin{subfigure}[t]{0.25\textwidth}
        \centering
        \includegraphics[width=\linewidth]{./assets/photos/reactor/reator-19-5.jpg}
        \caption{Murky suspension}
        
    \end{subfigure}\hspace{1em}
    \begin{subfigure}[t]{0.25\textwidth}
        \centering
        \includegraphics[width=\linewidth]{./assets/photos/reactor/reator-22-5.jpg}
        \caption{Black suspention}
        
    \end{subfigure}
    \caption{Visual changes in the reactor, within the first week.}
    
\end{figure}

At the end of the experimental period, after approximately two weeks of unattended operation, four discrete masses resembling fungal colonies were observed on the internal surfaces of the reactor during cleaning (?? and ??). The origin of these growths is uncertain, but they may have been introduced via the FW feed or during periods when the feeding tube was opened for repairs (e.g. leaks) and when changing the feed that the tube had to be cleaned. The presence of such organisms could have competed with the anaerobic microbial community for substrate, potentially affecting biogas production in the later stages of the experiment.

\begin{figure}[H]
    \centering
    \begin{subfigure}[t]{0.25\textwidth}
        \centering
        \includegraphics[width=\linewidth]{./assets/photos/reactor/reator-limp-seco-24-8.jpg}
        \caption{Masses inside the dry reactor}
        
    \end{subfigure}\hspace{1em}
    \begin{subfigure}[t]{0.25\textwidth}
        \centering
        \includegraphics[width=0.8\linewidth]{./assets/photos/reactor/reator-limp-agua-24-8.jpg}
        \caption{Masses floating in the reactor}
        
    \end{subfigure}
    \caption{Visual for the masses inside the reactor.}
    
\end{figure}

\section{Biogas Production}

Biogas production was monitored continuously using AMPTS II, which recorded both the cumulative volume of CH4 produced and the instantaneous flow rate. The entire experimental period is summarised in ?? and ??. To facilitate interpretation, the data have been divided into four distinct phases based on the feeding regime: (i) an adaptation phase~??, (ii) co-digestion of food waste and glycerol~??, (iii) digestion of crude glycerol~??, and (iv) a post-feeding period~??.

\subsection{Overall Production}

The cumulative CH4 production over the whole experiment (??) shows a continuous increase, indicating that the reactor maintained methanogenic activity throughout the monitored period. The total flow profile (??) exhibits distinct peaks corresponding to daily production of CH4, with variations in amplitude reflecting the different substrate compositions and organic loading rates applied during each phase. And a outlier peak at day 53, that was when the feed was changed from a mixture of FW and C-GL, to only C-GL, and due to some complications, created an anomaly in the graph.

\begin{figure}[H]
    \centering
    \VolumeGraphInt[xtick distance=10, minor xtick={1,2,...,97}]
        {data/processed/experimentos-day.csv}{0}{97}{Time (d)}{Volume (Nm L)}{}
    \caption{Cumulative volume of \ce{CH4} Nm L}
    
\end{figure}

\begin{figure}[H]
    \centering
    \FlowGraph[xtick distance=10, minor xtick={1,2,...,97}]
        {data/processed/experimentos-day.csv}{0}{97}{Time (d)}{Flow (Nm L/d)}{}
    \caption{Volume flow of \ce{CH4} Nm L/d}
    
\end{figure}

\subsection{Adaptation Phase}

During the initial adaptation phase, biogas production was low (?? and ??). This lag period is normal for anaerobic reactors as the microbial community adjusts to the new substrate and operating conditions. The small but measurable CH4 flow confirmed that the viability of the inoculum, and the gradual increase in flow suggested the onset of exponential growth of the microorganisms.

The spikes seen on the days 8 and 10 (??) probably happended by the first addition of C-GL to the feed.

\begin{figure}[H]
    \centering
    \VolumeGraphInt[xtick distance=10, minor xtick={1,2,...,97}]
        {data/processed/experimentos-day.csv}{0}{24}{Time (d)}{Volume (Nm L)}{}
    \caption{Cumulative volume of \ce{CH4} Nm L (Adaptation Phase)}
    
\end{figure}

\begin{figure}[H]
    \centering
    \FlowGraph[xtick distance=10, minor xtick={1,2,...,97}]
        {data/processed/experimentos-day.csv}{0}{24}{Time (d)}{Flow (Nm L/d)}{}
    \caption{Volume flow of \ce{CH4} Nm L/d (Adaptation Phase)}
    
\end{figure}

\subsection{Co-digestion of Food Waste and Crude Glycerol}

The reactor was fed with a mixture of FW and C-GL, a marked increase in biogas production was observed (?? and ??). The cumulative CH4 volume rose steeply, and the flow rate showed pronounced spikes immediately after each feeding event (??). These spikes are typical of readily degradable substrates, which are quickly hydrolysed and converted to VFA and subsequently to CH4. The average flow rate during this phase was higher than in any other phase, indicating that the co-digestion of FW and C-GL provided a balanced nutrient supply and avoided the inhibition often associated with high GL concentrations alone.

\begin{figure}[H]
    \centering
    \VolumeGraphInt[xtick distance=10, minor xtick={1,2,...,97}]
        {data/processed/experimentos-day.csv}{25}{52}{Time (d)}{Volume (Nm L)}{}
    \caption{Cumulative volume of \ce{CH4} Nm L (co-digestion of FW and C-GL)}
    
\end{figure}

\begin{figure}[H]
    \centering
    \FlowGraph[xtick distance=10, minor xtick={1,2,...,97}]
        {data/processed/experimentos-day.csv}{25}{52}{Time (d)}{Flow (Nm L/d)}{}
    \caption{Volume flow of \ce{CH4} Nm L/d (co-digestion of FW and C-GL)}
    
\end{figure}

\begin{figure}[H]
    \centering
    \FlowGraph[xtick distance=100, minor xtick={1,2,...,97}]
        {data/processed/experimentos-30min.csv}{600}{1248.75}{Time (h)}{Flow (Nm L/h)}{}
    \caption{Volume flow of \ce{CH4} Nm L/h (co-digestion of FW and C-GL)}
    
\end{figure}

\subsection{Digestion of Crude Glycerol}

When the feed was changed to C-GL alone, biogas production decreased (?? and ??). Although spikes in the flow were still observed after feeding (??), their magnitude was bigger than during the co-digestion phase, and the cumulative CH4 volume plateaued more quickly. Several factors may explain this lower performance:
\begin{itemize}
    \item The GL could have been transformed to fast or too slow, to maintaing a concentration of intermediate products.
    \item The retention rate of $\frac{1}{30}\,/d$ may have been insufficient to wash out inhibitory by-products or to maintain a stable population of glycerol fermenting bacteria.
    \item The appearance of fungal masses towards the end of this phase could have further stressed the anaerobic community by competing for substrate or introducing toxic metabolites. The fungal masses could have helped digest more complex sugars, so without a more thorough investigation of what it was, it can't be said it didn't help the AD system.
\end{itemize}

For those factors, can't be said that this phase is worse than the co-digestion phase. If the experiment was done in reverse, it's likely that we would have seen a reverse on the results, were the phase with only C-GL would have a bigger production than the co-digestion of FW and C-GL.

\begin{figure}[H]
    \centering
    \VolumeGraphInt[xtick distance=10, minor xtick={1,2,...,97}]
        {data/processed/experimentos-day.csv}{53}{81}{Time (d)}{Volume (Nm L)}{}
    \caption{Cumulative volume of \ce{CH4} Nm L (Anaerobic digestion of C-GL)}
    
\end{figure}

\begin{figure}[H]
    \centering
    \FlowGraph[xtick distance=10, minor xtick={1,2,...,97}]
        {data/processed/experimentos-day.csv}{53}{81}{Time (d)}{Flow (Nm L/d)}{}
    \caption{Volume flow of \ce{CH4} Nm L/d (Anaerobic digestion of C-GL)}
    
\end{figure}

\begin{figure}[H]
    \centering
    \FlowGraph[xtick distance=100, minor xtick={1,2,...,97}]
        {data/processed/experimentos-30min.csv}{1272}{1941.5}{Time (h)}{Flow (Nm L/h)}{}
    \caption{Volume flow of \ce{CH4} Nm L/h (Anaerobic digestion of C-GL)}
    
\end{figure}

\subsection{Post-feeding Period}

During the final two weeks, no fresh substrate was added to the reactor, it just had some residual substrate from the C-GL only phase for a week. Nevertheless, a small residual biogas production was recorded (?? and ??). This production can be attributed to the slow degradation of accumulated organic matter and microbial decay products. The flow rate declined gradually, confirming that the reactor was no longer being actively fed and that the remaining substrate was limited.

\begin{figure}[H]
    \centering
    \VolumeGraphInt[xtick distance=10, minor xtick={1,2,...,97}]
        {data/processed/experimentos-day.csv}{82}{97}{Time (d)}{Volume (Nm L)}{}
    \caption{Cumulative volume of \ce{CH4} Nm L (Post-feeding period)}
    
\end{figure}

\begin{figure}[H]
    \centering
    \FlowGraph[xtick distance=10, minor xtick={1,2,...,97}]
        {data/processed/experimentos-day.csv}{82}{97}{Time (d)}{Flow (Nm L/d)}{}
    \caption{Volume flow of \ce{CH4} Nm L/d (Post-feeding period)}
    
\end{figure}

\subsection{Alkalinity and Volatile Fatty Acids}

The total Alk and VFA concentrations were monitored throughout the experiment to assess process stability. The results are shown in ??. Initially, both Alk and VFA were low, reflecting the early adaptation phase before buffer addition. After the manual introduction of KHCO3, Alk increased sharply to values above 7000 mg Ca CO3/L and remained relatively high until close to the end. This high Alk was maintained by the addition of bicarbonate to the feed, which helped to counteract the acidity produced during the degradation of the feed.

Volatile fatty acid concentrations also increased after the start of feeding, reaching values between 3000–5000 mg CH3COOH/L. Such high VFA levels are often associated with process imbalance, as they can inhibit methanogenesis. However, the high Alk provided a strong buffering capacity, and the pH remained above 6.5 for most of the experiment, with the co-digestion the pH was falling and with C-GL only the pH was increasing. The vfa/alk is a more reliable indicator of instability. In this experiment, the ratio fluctuated between 0.3–1.2. While using the co-digestion feed a ration around 0.8–1.0 had the highest production.

\begin{figure}[H]
    \centering
    \begingroup
    \pgfplotstableread[col sep=comma]{data/processed/analysis.csv}{\datatable}
    \resizebox{\textwidth}{!}{\begin{tikzpicture}
\begin{axis}[
            width=15cm, height=10cm,
            xlabel={Time (d)},
            ylabel={Alk (mg Ca CO3/L)},
            xmin=0, ymin=0,
            axis y line*=left,        legend pos=south east,
        ]
            \addplot[
                color=blue, thick, mark=*, mark size=1.5pt,
            ] table[x index=0, y index=2] {\datatable};
            \addlegendentry{Alk}

\addlegendimage{color=red, thick, mark=square*, mark size=1.5pt}
            \addlegendentry{VFA}
        \end{axis}

\begin{axis}[
            width=15cm, height=10cm,
            xmin=0, ymin=0,
            axis y line*=right,       axis x line=none,         ylabel={VFA (mg CH3COOH/L)},
]
            \addplot[
                color=red, thick, mark=square*, mark size=1.5pt,
            ] table[x index=0, y index=3] {\datatable};
        \end{axis}
    \end{tikzpicture}}
    \endgroup
    \caption{Alkalinity \& volatile fatty acid comparison graph}
    
\end{figure}

\begin{figure}[H]
    \centering
    \begingroup
    \pgfplotstableread[col sep=comma]{data/processed/analysis.csv}{\datatable}
    \resizebox{\textwidth}{!}{\begin{tikzpicture}
\begin{axis}[
            width=15cm, height=10cm,
            xlabel={Time (d)},
            ylabel={pH},
            xmin=0, ymin=0,
            legend style={
                at={(0.98, 0.4)},     anchor=east,          draw=black,
                fill=white,
                fill opacity=0.9,
            },
            ]
            \addplot[
                color=green, thick, mark=*, mark size=1.5pt,
            ] table[x index=0, y index=1] {\datatable};
            \addlegendentry{pH}

\addlegendimage{color=purple, thick, mark=square*, mark size=1.5pt}
            \addlegendentry{vfa/alk}
        \end{axis}

\begin{axis}[
            width=15cm, height=10cm,
            xmin=0, ymin=0,
            axis y line*=right,       axis x line=none,         ylabel={vfa/alk},
]
            \addplot[
                color=purple, thick, mark=square*, mark size=1.5pt,
            ] table[x index=0, y index=4] {\datatable};
        \end{axis}
    \end{tikzpicture}}
    \endgroup
    \caption{pH \& vfa/alk comparison graph}
    
\end{figure}

\clearpage{}
\clearpage{}\chapter{Conclusion}

A direct comparison of the co-digestion and C-GL only phases reveals that the co-digestion of FW and C-GL resulted in higher biogas yields and more stable operation. The pH was decreasing faster with the co-digestion, where with C-GL only phase it wasn't decreasing that fast, and was even growing with the KHCO3 that was added to maintain the pH. These results support the hypothesis that C-GL, when co-digested with FW, can enhance biogas production by balancing the C/N and providing readily available carbon, whereas C-GL alone may work, but would need a different setup and ajustments.

The C-GL only phase, can't be said to be strictly worse than the co-digestion of FW and C-GL, since both phases were done in sequence in the same reactor, this could have given the C-GL pahse disavantages on a contaminated enviroment.

Using C-GL, with some P-GL to make the necessary ajustment to balance the C/N, could be a viable solution for the C-GL problem. It's only major problem, is that C-GL contains a lot of inpurities that can affect the anaerobic digestion, and for a system that can reliably consume it, it will need a better control of it's substrate.

During the whole process, the biggest production of CH4 was when vfa/alk was \approx1 with a pH of \approx6.5, during the C-GL only phase, even with not optmal valuos of pH and vfa/alk, it was able to maintain a resonable amount of production of CH4, showing that it can be a suitable substrate for AD.

Future studies could involve on how to use better the C-GL as a feed, to try subside it's downside of it's impurities and need for an extra caron rich substrate to balance the feed.
\clearpage{}

\printbibliography[heading=bibintoc, title={References}]

\end{document}
